\section*{第 \ref{ch:g10:vectors} 章 --- 向量}

\begin{solution}{exo:g10:vectors:1}
$\vect{AB} + \vect{BD} = \vect{AD}$（沙勒关系）。

$\vect{KL} + \vect{LM} + \vect{MK} = \vect{KK} = \vec 0$.

$\vect{AC} - \vect{BC} = \vect{AC} + \vect{CB} = \vect{AB}$.

$\vect{AB} + \vect{CA} = \vect{CA} + \vect{AB} = \vect{CB}$.
\end{solution}

\begin{solution}{exo:g10:vectors:2}
$\vect{AB}\,(5-2,\ 3-1) = (3, 2)$; \quad
$\vect{BC}\,(-1-5,\ 4-3) = (-6, 1)$; \quad
$\vect{AC}\,(-1-2,\ 4-1) = (-3, 3)$.
验算：$(3, 2) + (-6, 1) = (-3, 3)$，按坐标逐一相加。
\end{solution}

\begin{solution}{exo:g10:vectors:3}
$\vec u + \vec v = (2 + (-1),\ -3 + 4) = (1, 1)$.

$3\vec u = (6, -9)$.

$2\vec u - \vec v = (4 - (-1),\ -6 - 4) = (5, -10)$.

$\norm{\vec u} = \sqrt{2^2 + (-3)^2} = \sqrt{13}$.
\end{solution}

\begin{solution}{exo:g10:vectors:4}
(a) $4 \times 9 - 6 \times 6 = 36 - 36 = 0$：共线
($\vec v = \frac32 \vec u$).

(b) $2 \times 10 - (-5) \times (-4) = 20 - 20 = 0$：共线
($\vec v = -2\vec u$).

(c) $3 \times 3 - 1 \times 1 = 8 \neq 0$：不共线。
\end{solution}

\begin{solution}{exo:g10:vectors:5}
$\vect{AB}\,(4, -2)$，$\vect{DC}\,(6 - x_D,\ 3 - y_D)$。条件
$\vect{AB} = \vect{DC}$ 给出 $6 - x_D = 4$ 且 $3 - y_D = -2$，于是
$D(2, 5)$。
\end{solution}

\begin{solution}{exo:g10:vectors:6}
\emph{1.} $\vect{AB}\,(3, 2)$ 与 $\vect{AC}\,(6, -2)$，于是
$\vect{AB} + \vect{AC} = (9, 0)$，且
$M = (-2 + 9,\ 1 + 0) = (7, 1)$.

\emph{2.} 由平行四边形法则，$\vect{AM} = \vect{AB} + \vect{AC}$
表示 $ABMC$ 是平行四边形（$M$ 为与
$A$ 相对的第四个顶点）。可以验证：$\vect{AB}\,(3,2)$ 与 $\vect{CM}\,(7-4,\ 1-(-1)) =
(3, 2)$ 相等。
\end{solution}

\begin{solution}{exo:g10:vectors:7}
第一组：$\vect{AB}\,(3, 2)$，$\vect{AC}\,(9, 6)$；
$3 \times 6 - 2 \times 9 = 0$：共线（$\vect{AC} = 3\vect{AB}$）。

第二组：$\vect{AB}\,(2, -1)$，$\vect{AC}\,(8, -4)$；
$2 \times (-4) - (-1) \times 8 = -8 + 8 = 0$：同样共线
($\vect{AC} = 4\vect{AB}$).
\end{solution}

\begin{solution}{exo:g10:vectors:8}
$\vect{AB}\,(4, 2)$ 与 $\vect{DC}\,(6-2,\ 1-(-1)) = (4, 2)$：两个
向量\emph{相等}，因此 $ABCD$ 是平行四边形，特别地
$(AB) \parallel (CD)$。平行四边形的性质是较强的那个：它
蕴含平行（即向量共线），反之不成立。
\end{solution}

\begin{solution}{exo:g10:vectors:9}
\emph{1.} $\vect{AB}\,(6, 3)$，于是
$\vect{AI} = \frac23\vect{AB} = (4, 2)$，故 $I(4, 2)$。

\emph{2.} $\vect{CB}\,(6-2,\ 3-5) = (4, -2)$，于是
$\vect{CJ} = \frac23\vect{CB} = \left(\frac83, -\frac43\right)$，故
$J\left(2 + \frac83,\ 5 - \frac43\right) = \left(\frac{14}{3},
\frac{11}{3}\right)$.

\emph{3.} $\vect{IJ}\,\left(\frac{14}{3} - 4,\ \frac{11}{3} - 2\right) =
\left(\frac23, \frac53\right)$ 与 $\vect{AC}\,(2, 5)$。共线判别
准则：
\[
\frac23 \times 5 - \frac53 \times 2 = \frac{10}{3} - \frac{10}{3} = 0 :
\]
向量共线，故 $(IJ)$ 与 $(AC)$ 平行。（$I$ 与
$J$ 都在所在边朝 $B$ 方向的三分之一处：小
三角形 $IBJ$ 是 $ABC$ 的缩小图形。）
\end{solution}

\begin{solution}{exo:g10:vectors:10}
各中点分别为
\[
I\left(\tfrac{x_A + x_B}{2}, \tfrac{y_A + y_B}{2}\right), \quad
J\left(\tfrac{x_B + x_C}{2}, \tfrac{y_B + y_C}{2}\right), \quad
K\left(\tfrac{x_C + x_D}{2}, \tfrac{y_C + y_D}{2}\right), \quad
L\left(\tfrac{x_D + x_A}{2}, \tfrac{y_D + y_A}{2}\right).
\]
两个向量的横坐标：
\[
x_J - x_I = \frac{x_B + x_C}{2} - \frac{x_A + x_B}{2}
= \frac{x_C - x_A}{2},
\qquad
x_K - x_L = \frac{x_C + x_D}{2} - \frac{x_D + x_A}{2}
= \frac{x_C - x_A}{2},
\]
对纵坐标同理可得。于是
$\vect{IJ} = \vect{LK}$（二者均等于 $\frac12\vect{AC}$），且 $IJKL$ 是
平行四边形，无论四边形 $ABCD$ 形状如何。
\end{solution}

\begin{solution}{pb:g10:vectors:1}
\textbf{1.} $\vect{AB} + \vect{BC} + \vect{CD} = \vect{AD}$;
\quad $\vect{MA} - \vect{MB} = \vect{MA} + \vect{BM} =
\vect{BA}$; \quad $\vect{AB} + \vect{BC} + \vect{CD} +
\vect{DA} = \vect{AA} = \vec 0$（调整第三个和式的次序）。

\textbf{2.} $\vect{OB} - \vect{OA} = \vect{AO} + \vect{OB} =
\vect{AB}$（由沙勒关系）。

\textbf{3.} 当且仅当 $I$ 为 $[AB]$ 的中点时
$\vect{IA} = -\vect{IB}$，即\
$\vect{IA} + \vect{IB} = \vec 0$。于是对任意 $O$：
$\vec 0 = \vect{IO} + \vect{OA} + \vect{IO} + \vect{OB}$，故
$2\,\vect{OI} = \vect{OA} + \vect{OB}$ --- 反之亦然。

\textbf{4.} $\vect{AB} = (4, 2)$，且 $D$ 满足
$\vect{DC} = \vect{AB}$：$D = C - (4, 2) = (2, -2)$ --- 与
\cref{pb:g10:coordgeom:1}中的 $D$ 相同，而此法更为便捷。

\textbf{5.} 向量 $\vec u$ 的平移 --- 恰是
由两面平行镜子或两次
中心对称所产生的变换。$(1, 2)$ 的像为 $(4, 1)$。

\textbf{6.} 将每个 $\vect{GX}$ 展开为
$\vect{GO} + \vect{OX}$:
$3\,\vect{GO} + \vect{OA} + \vect{OB} + \vect{OC} = \vec 0$,
故 $\vect{OG} = \frac13(\vect{OA} + \vect{OB} + \vect{OC})$：
满足条件的点恰有一个，
其坐标为三个顶点坐标的平均值。

\textbf{7.} $\vect{AB} + \vect{AC} =
(\vect{AG} + \vect{GB}) + (\vect{AG} + \vect{GC})
= 2\vect{AG} + (\vect{GB} + \vect{GC})
= 2\vect{AG} - \vect{GA} = 3\vect{AG}$，由此
$\vect{AG} = \frac13(\vect{AB} + \vect{AC})$。中点
公式（第 3 题，从 $A$ 出发）给出
$\vect{AK} = \frac12(\vect{AB} + \vect{AC})$。比较可得：
$\vect{AG} = \frac23\vect{AK}$ --- $G$ 位于中线
$[AK]$ 上，位于距 $A$ 三分之二处。

\textbf{8.} 定义条件对 $A$、$B$、$C$
一视同仁，故从 $B$ 或 $C$ 出发的同样计算，
会将\emph{同一个} $G$ 置于这些中线的三分之二处：
共点并不需额外代价。平行四边形的证法
虽巧妙却属特设；坐标法可以计算，却不够清晰；
向量法使对称性一目了然，证明只需两行。

\textbf{9.} $G\left(\frac{1+5+3}{3}, \frac{1+2+6}{3}\right) =
(3, 3)$；$K(4, 4)$；$\vect{AG} = (2, 2)$ 及
$\vect{AK} = (3, 3)$：行列式
$2 \times 3 - 2 \times 3 = 0$，共线，
且事实上 $\vect{AG} = \frac23\vect{AK}$。

\textbf{10.} $G'\left(\frac{2 + 5 + 3}{4}, \frac{2 + 2 +
6}{4}\right) = (2.5,\ 2.5)$：由 $G(3,3)$ 拉向
权重加倍的顶点 $A(1,1)$，正如平衡点应有的那样。

\textbf{11.} $3 \times 4 - (-2)(-6) = 12 - 12 = 0$：共线
($\vec v = -2\vec u$). $2 \times 9 - 5 \times 4 = -2 \neq 0$:
不共线。

\textbf{12.} $\vect{AB} = (2, 3)$, $\vect{AC} = (6, 9)$:
行列式 $2 \times 9 - 3 \times 6 = 0$：共线。

\textbf{13.} $\vect{OI} = \frac12(\vect{OA} + \vect{OB})$，而
$\vect{OK} = \frac12(\vect{OC} + \vect{OD})$，故
$[IK]$ 的中点的位置向量为
$\frac12(\vect{OI} + \vect{OK}) = \frac14(\vect{OA} +
\vect{OB} + \vect{OC} + \vect{OD})$ --- 而对 $[JL]$ 的计算
给出的正是同一向量。两条双中线（以及
由瓦里尼翁定理，平行四边形 $IJKL$ 的对角线）都经过
这一点，即四边形的重心。

\textbf{14.} $\vect{MN} = \vect{AN} - \vect{AM} =
\frac13\vect{AC} - \frac13\vect{AB} = \frac13\vect{BC}$:
故 $(MN) \parallel (BC)$，且 $MN = \frac13 BC$。

\textbf{15.} 设 $\vect{AM} = k\,\vect{AB}$，且
$\vect{AN} = k\,\vect{AC}$：$\vect{MN} = k\,\vect{BC}$ --- 一
行中既包含中点定理（$k = \frac12$），
也包含泰勒斯定理（任意 $k$）：比值在做减法时
丝毫未受影响。

\textbf{16.} 合速度 $(3, 4)$，大小
$\sqrt{9 + 16} = 5$ km/h：小船实际
东北偏北方向，沿 $3$--$4$--$5$ 三角形的
斜边行进。

\textbf{17.} $\vec u + \vec v + \vec w = (0, 0)$：圆环
保持静止 --- 即平衡。沿闭合多边形
$A_1 A_2 \dots A_n A_1$，由沙勒关系，和式逐项相消为
$\vect{A_1 A_1} = \vec 0$：绕行一周即回到出发点。

\textbf{18.} $\vec w = -(\vec u + \vec v) = -(3, 2) =
(-3, -2)$.

\textbf{19.} 合速度 $(1.5,\ 2)$，大小
$\sqrt{2.25 + 4} = 2.5$ m/s。渡河：$50$ m 宽的河，以 $2$ m/s 横渡
需 $25$ s；漂移：$1.5 \times 25 = 37.5$ m，向下游。

\textbf{20.} 本周末重新证明了以下内容：重心的
三分之二定理（第 7 题）、双中线与瓦里尼翁中心（第
13 题）、泰勒斯与中点定理（第 14--15 题）。
箭头仍做不到的是：测量\emph{角}（从而测量
任意方向上的长度） --- 十一年级学习的数量积
恰是所缺的那件工具。
\end{solution}
